\section{Conclusion}

Three statements, in the order they should be made.

Optimal saving rises with cash on hand for every continuation value, so a value function made
non-concave by a means test cannot by itself reverse a policy. The common argument from
non-concavity to non-monotonicity is therefore not valid as stated, and it is not a matter of
sharper hypotheses: the continuation cancels.

Consumption is a different object and it is genuinely non-monotone. At the wealth where the
household stops protecting its pension, saving jumps up by more than wealth has risen and
consumption drops. We give the drop exactly in a two-period example and measure it in a calibrated
retiree, where Australia's asset test produces a fall of two and a half to five per cent of mean
earnings at the free area, at every age of retirement and in the last year of work as well.

Next period's assets rise with this period's, at each earnings shock, again for every continuation.
This fails in one case only: when the test is assessed on the assets currently held and withdraws
the pension faster than assets earn, cash on hand itself falls and the policy follows. The
threshold is the taper rate against the gross return, and it is sharp.

For applied work the practical content is that a means test of the Australian kind leaves the
saving policy monotone and puts a discontinuity in the consumption function at the free area.
Solution methods that rely on monotonicity of the policy in assets are safe; methods that rely on a
continuous or monotone consumption function, or on concavity of the value function, are not. That
is close to the practice the applied literature has already settled on, which abandons Euler
equation methods and searches a fine grid \citep{Wellschmied2021,GarstenauerSiassi2024}, and the
results here say which half of that caution is needed. For
theory the content is that the two state spaces are not interchangeable once the pension depends on
assets, and that monotone comparative statics survives non-concavity far better than it is usually
given credit for.
